\thusetup{
  %******************************
  % 注意：
  %   1. 配置里面不要出现空行
  %   2. 不需要的配置信息可以删除
  %******************************
  %
  %=====
  % 秘级
  %=====
  % secretlevel={秘密},
  % secretyear={10},
  %
  %=========
  % 中文信息
  %=========
  ctitle={基于BLE Mesh的智能家居平台\\研究报告},
  cdepartment={金苹果锦城第一中学},
  cmajor={计算机科学与技术},
  cauthor={李董睿煊、翟羽佳},
  csupervisor={伍冬莉},
  % 日期自动使用当前时间，若需指定按如下方式修改：
  % cdate={超新星纪元},
  %
  % 博士后专有部分
  % catalognumber     = {分类号},  % 可以留空
  % udc               = {UDC},  % 可以留空
  % id                = {编号},  % 可以留空： id={},
  % cfirstdiscipline  = {计算机科学与技术},  % 流动站（一级学科）名称
  % cseconddiscipline = {系统结构},  % 专 业（二级学科）名称
  % postdoctordate    = {2009 年 7 月——2011 年 7 月},  % 工作完成日期
  % postdocstartdate  = {2009 年 7 月 1 日},  % 研究工作起始时间
  % postdocenddate    = {2011 年 7 月 1 日},  % 研究工作期满时间
  %
  %=========
  % 英文信息
  %=========
  % etitle={An Introduction to \LaTeX{} Thesis Template of Tsinghua University v\version},
  % 这块比较复杂，需要分情况讨论：
  % 1. 学术型硕士
  %    edegree：必须为Master of Arts或Master of Science（注意大小写）
  %             “哲学、文学、历史学、法学、教育学、艺术学门类，公共管理学科
  %              填写Master of Arts，其它填写Master of Science”
  %    emajor：“获得一级学科授权的学科填写一级学科名称，其它填写二级学科名称”
  % 2. 专业型硕士
  %    edegree：“填写专业学位英文名称全称”
  %    emajor：“工程硕士填写工程领域，其它专业学位不填写此项”
  % 3. 学术型博士
  %    edegree：Doctor of Philosophy（注意大小写）
  %    emajor：“获得一级学科授权的学科填写一级学科名称，其它填写二级学科名称”
  % 4. 专业型博士
  %    edegree：“填写专业学位英文名称全称”
  %    emajor：不填写此项
  % edegree={Doctor of Engineering},
  % emajor={Computer Science and Technology},
  % eauthor={Xue Ruini},
  % esupervisor={Professor Zheng Weimin},
  % eassosupervisor={Chen Wenguang},
  % 日期自动生成，若需指定按如下方式修改：
  % edate={December, 2005}
  %
  % 关键词用“英文逗号”分割
  ckeywords={TeX, LaTeX, CJK, 模板, 论文},
  % ekeywords={TeX, LaTeX, CJK, template, thesis}
}
\begin{schoollogo}
  \begin{figure}[H] % use float package if you want it here
    \centering
    \includegraphics{logo.png}
  \end{figure}
\end{schoollogo}
% 定义中英文摘要和关键字
\begin{cabstract}
  论文的摘要是对论文研究内容和成果的高度概括。摘要应对论文所研究的问题及其研究目
  的进行描述，对研究方法和过程进行简单介绍，对研究成果和所得结论进行概括。摘要应
  具有独立性和自明性，其内容应包含与论文全文同等量的主要信息。使读者即使不阅读全
  文，通过摘要就能了解论文的总体内容和主要成果。

  论文摘要的书写应力求精确、简明。切忌写成对论文书写内容进行提要的形式，尤其要避
  免“第 1 章……；第 2 章……；……”这种或类似的test。

  本文介绍清华大学论文模板 \thuthesis{} 的使用方法。本模板符合学校的本科、硕士、
  博士论文格式要求。

  本文的创新点主要有：
  \begin{itemize}
    \item 用例子来解释模板的使用方法；
    \item 用废话来填充无关紧要的部分；
    \item 一边学习摸索一边编写新代码。
    \item 巴拉巴拉
  \end{itemize}

  关键词是为了文献标引工作、用以表示全文主要内容信息的单词或术语。关键词不超过 5
  个，每个关键词中间用分号分隔。（模板作者注：关键词分隔符不用考虑，模板会自动处
  理。英文关键词同理。）
\end{cabstract}

% 如果习惯关键字跟在摘要文字后面，可以用直接命令来设置，如下：
% \ckeywords{\TeX, \LaTeX, CJK, 模板, 论文}

% \begin{eabstract}
%    An abstract of a dissertation is a summary and extraction of research work
%    and contributions. Included in an abstract should be description of research
%    topic and research objective, brief introduction to methodology and research
%    process, and summarization of conclusion and contributions of the
%    research. An abstract should be characterized by independence and clarity and
%    carry identical information with the dissertation. It should be such that the
%    general idea and major contributions of the dissertation are conveyed without
%    reading the dissertation.

%    An abstract should be concise and to the point. It is a misunderstanding to
%    make an abstract an outline of the dissertation and words ``the first
%    chapter'', ``the second chapter'' and the like should be avoided in the
%    abstract.

%    Key words are terms used in a dissertation for indexing, reflecting core
%    information of the dissertation. An abstract may contain a maximum of 5 key
%    words, with semi-colons used in between to separate one another.
% \end{eabstract}

% \ekeywords{\TeX, \LaTeX, CJK, template, thesis}
